\section*{Capítulo \ref{ch:b2:population-genetics} --- Genética de populações}

\begin{solution}{exo:b2:population-genetics:1}
$p(M) = (1280 + 320)/2000 = 0.8$, $q(N) = 0.2$. Esperados $MM$:
$0.64\times 1000 = 640$; $MN$: $2\times 0.8\times 0.2\times 1000 =
320$; $NN$: 40 --- exatamente os números observados. A população está
em equilíbrio nesse lócus.
\end{solution}

\begin{solution}{exo:b2:population-genetics:2}
$q = \sqrt{1/20000} = 0.0071$; portadores $2pq \approx 0.014$, uma
pessoa em 70. Os portadores superam os afetados na razão de $2p/q \approx
280$ para um.
\end{solution}

\begin{solution}{exo:b2:population-genetics:3}
\omterm{def:b2:population-genetics:fitness}{Aptidão}: o número esperado de descendentes de um \omterm{def:b2:meiosis-heredity:vocabulary}{genótipo} em relação ao
melhor. \omterm{def:b2:population-genetics:fitness}{Coeficiente de seleção} $s$: o déficit de \omterm{def:b2:population-genetics:fitness}{aptidão}, $w = 1 -
s$. Herdabilidade $h^{2}$: a fração da variância fenotípica que é
genética aditiva, ou seja, a inclinação da regressão dos descendentes
sobre os pais. Forças: \omterm{def:b2:genome-diversification:mutation}{mutação}, seleção, deriva, migração, cruzamentos
não aleatórios (a última muda as frequências genotípicas, e não as
alélicas).
\end{solution}

\begin{solution}{exo:b2:population-genetics:4}
Direcional: o melanismo da mariposa-salpicada, tendo como agente as
aves que caçam pela visão sobre troncos escurecidos pela fuligem.
Estabilizadora: o peso ao nascer humano, tendo como agente a mortalidade
dos bebês muito pequenos e muito grandes. Balanceadora: a anemia
falciforme, tendo como agentes a malária (contra $AA$) e a anemia (contra
$SS$) em conjunto.
\end{solution}

\begin{solution}{exo:b2:population-genetics:5}
$q_t = q_0/(1 + tq_0)$: reduzir à metade leva $1/q_0 = 50$ gerações;
chegar a $0.001$ leva $1/0.001 - 1/0.02 = 950$ gerações. Quase todas as
cópias do \omterm{def:b2:meiosis-heredity:vocabulary}{alelo} estão em \omterm{def:b2:meiosis-heredity:vocabulary}{heterozigotos}, invisíveis para a seleção;
impedir que os afetados se reproduzam (o que a natureza já faz, pois eles
morrem) não muda nada de mensurável ao longo de uma vida humana.
\end{solution}

\begin{solution}{exo:b2:population-genetics:6}
$p = 0.3$: $\bar w = 0.09 + 0.42 + 0.8\times 0.49 = 0.902$; $p' =
(0.09 + 0.21)/0.902 = 0.333$; $\Delta p = +0.033$. $p = 0.9$: $\bar w
= 0.998$, $p' = 0.9018$, $\Delta p = +0.0018$. Com $p = 0.9$, apenas
$q^{2} = 0.01$ da população é $aa$: o \omterm{def:b2:meiosis-heredity:vocabulary}{alelo} contra o qual a seleção age
está escondido nos \omterm{def:b2:meiosis-heredity:vocabulary}{heterozigotos}, e $\Delta p = spq^{2}/\bar w$
carrega o fator $q^{2}$.
\end{solution}

\begin{solution}{exo:b2:population-genetics:7}
Com $s = 0.12$ contra $AA$ e $t = 0.8$ contra $SS$, $\hat q =
s/(s + t) = 0.12/0.92 = 0.13$. Nascimentos afetados $\hat q^{2} = 0.017$,
cerca de uma criança em 60 --- o preço da proteção de que gozam os
\omterm{def:b2:meiosis-heredity:vocabulary}{heterozigotos}.
\end{solution}

\begin{solution}{exo:b2:population-genetics:8}
Letal \omterm{def:b2:meiosis-heredity:vocabulary}{recessivo}: $\hat q = \sqrt{\mu} = \sqrt{2\times 10^{-5}} =
0.0045$; nascimentos afetados $\hat q^{2} = \mu = 2\times 10^{-5}$, um em
\num{50000}. Letal \omterm{def:b2:meiosis-heredity:vocabulary}{dominante}: toda cópia é eliminada na geração em que
aparece, de modo que os nascimentos afetados são simplesmente as novas
\omterm{def:b2:genome-diversification:mutation}{mutações}, $2\mu = 4\times 10^{-5}$ (dois gametas por nascimento), um em
\num{25000} --- o \omterm{def:b2:meiosis-heredity:vocabulary}{recessivo} esconde seus \omterm{def:b2:meiosis-heredity:vocabulary}{alelos} na razão de duzentos
para um, o \omterm{def:b2:meiosis-heredity:vocabulary}{dominante} não esconde nada.
\end{solution}

\begin{solution}{exo:b2:population-genetics:9}
$H_t = 0.5\,(1 - 1/100)^{t}$: após 10 gerações, $0.45$; após 50,
$0.30$; após 200, $0.067$. Para $(1 - 1/2N)^{100} = 0.9$:
$100/(2N) \approx -\ln 0.9 = 0.105$, logo $N \approx 475$, cerca de 500
indivíduos reprodutores.
\end{solution}

\begin{solution}{exo:b2:population-genetics:10}
Vinte cópias gênicas: $P(\text{ausente}) = 0.9^{20} = 0.12$. Uma única
cópia entre 20 tem frequência $0.05$, e um \omterm{def:b2:meiosis-heredity:vocabulary}{alelo} neutro se fixa com
probabilidade igual à sua frequência: $0.05$. Num continente de, digamos,
\num{100000} aves, uma única cópia tem probabilidade $5\times 10^{-6}$
de se fixar. A ilha é um lugar onde os \omterm{def:b2:meiosis-heredity:vocabulary}{alelos} raros se perdem e onde os
que sobrevivem podem dominar: a deriva ao mesmo tempo empobrece e
diferencia.
\end{solution}

\begin{solution}{exo:b2:population-genetics:11}
$R = h^{2}S = 0.3\times 1.4\times \qty{1000}{L} = \qty{420}{L}$ por
geração. Com touros do \qty{1}{\%} melhor, o diferencial médio é de
$(1.4 + 2.7)/2 = 2.05$ desvios-padrão e $R =
0.3\times 2050 = \qty{615}{L}$. A resposta desacelera porque a seleção
consome a variância aditiva --- os \omterm{def:b2:meiosis-heredity:vocabulary}{alelos} favoráveis chegam à fixação e
$h^{2}$ cai --- e porque as mudanças correlacionadas em outros traços
(fertilidade, saúde) impõem custos que a seleção apenas sobre a produção
não enxerga.
\end{solution}

\begin{solution}{exo:b2:population-genetics:12}
A seleção enxerga \omterm{def:b2:meiosis-heredity:vocabulary}{fenótipos}. Um \omterm{def:b2:meiosis-heredity:vocabulary}{alelo} \omterm{def:b2:meiosis-heredity:vocabulary}{recessivo} num \omterm{def:b2:meiosis-heredity:vocabulary}{heterozigoto} não
produz \omterm{def:b2:meiosis-heredity:vocabulary}{fenótipo}, de modo que a seleção não consegue eliminá-lo, e sua
frequência cai como $1/t$, cada vez mais devagar. Uma \omterm{def:b2:genome-diversification:mutation}{mutação} neutra não
produz diferença de \omterm{def:b2:population-genetics:fitness}{aptidão}; seu destino depende só da deriva, e ela se
fixa com probabilidade $1/2N$. Um caráter de herdabilidade nula pode ser
fortemente selecionado --- os sobreviventes podem diferir muito da
população ---, e no entanto a geração seguinte fica inalterada, porque as
diferenças não eram herdadas. A seleção só pode agir sobre diferenças
herdáveis que fazem diferença para a \omterm{def:b2:population-genetics:fitness}{aptidão}; todo o resto evolui por
acaso ou não evolui.
\end{solution}

\begin{solution}{pb:b2:population-genetics:1}
\textbf{1.} Fração melânica $1 - q^{2} = 1 - 0.999^{2} = 0.002$,
uma mariposa em 500; a fração dos \omterm{def:b2:meiosis-heredity:vocabulary}{alelos} $M$ em \omterm{def:b2:meiosis-heredity:vocabulary}{heterozigotos} é
$2pq/(2pq + 2p^{2}) = q = 0.999$: praticamente todos.
\textbf{2.} Com $\bar w = 1 - sq^{2}$, as cópias do \omterm{def:b2:meiosis-heredity:vocabulary}{alelo} claro
após a seleção são $pq\cdot 1 + q^{2}(1 - s)$ de um total $\bar w$, logo $q'
= (pq + q^{2} - sq^{2})/(1 - sq^{2}) = q(1 - sq)/(1 - sq^{2})$.
\textbf{3.} $q_1 = 0.999\times(1 - 0.2997)/(1 - 0.2994) = 0.99857$:
uma mudança de $-0.0014$.
\textbf{4.} $q = \sqrt{0.1} = 0.32$.
\textbf{5.} A recorrência, iterada, leva 28 gerações para $s =
0.3$; a aproximação dá a mesma ordem (a mudança é de
$-0.0003$ no início e depois acelera à medida que $p$ cresce). As 50
gerações observadas correspondem a um $s$ menor, cerca de $0.2$ (a recorrência
dá 45 gerações para $s = 0.2$): uma seleção de \qtyrange{20}{30}{\%}
por geração, enorme para os padrões da maioria dos \omterm{def:b2:meiosis-heredity:vocabulary}{alelos}.
\textbf{6.} Com $w_{MM} = w_{Mm} = 1 - s$, $w_{mm} = 1$: $\bar w = 1
- s(1 - q^{2})$ e $p' = p(1 - s)/\bar w$, logo $\Delta p = p[(1 - s) -
\bar w]/\bar w = -spq^{2}/\bar w$. A forma clara só é favorecida como
o \omterm{def:b2:meiosis-heredity:vocabulary}{homozigoto} $mm$, uma fração $q^{2} = 0.1$ no início: o
retorno começa devagar. Quando $q \to 1$, $\Delta p \approx -sp$ e $M$
declina exponencialmente, por um fator $1 - s$ a cada geração: um
\omterm{def:b2:meiosis-heredity:vocabulary}{alelo} \omterm{def:b2:meiosis-heredity:vocabulary}{dominante} não tem onde se esconder. A ascensão foi a imagem
espelhada: $\Delta p \approx +sp$ enquanto $M$ era raro (início
exponencial), depois
$\Delta q \approx -sq^{2}p \to$ um arrastar à medida que o \omterm{def:b2:meiosis-heredity:vocabulary}{alelo} claro se
escondia nos \omterm{def:b2:meiosis-heredity:vocabulary}{heterozigotos}. A recorrência dá 27 gerações para o retorno
de $p = 0.68$ a $0.001$, 16 delas para levar o \omterm{def:b2:meiosis-heredity:vocabulary}{alelo} a \qty{5}{\%}.
\textbf{7.} $24$ cópias gênicas: $0.95^{24} = 0.29$.
\textbf{8.} $H_1 = 0.4\,(1 - 1/24) = 0.383$.
\textbf{9.} $H_{100} = 0.383\,(1 - 1/80)^{100} = 0.383\times 0.284 =
0.109$: \qty{27}{\%} dos $0.4$ do continente.
\textbf{10.} $P(\text{fixação}) = 1/(2N) = 1/80 = 0.0125$; tempo
esperado até a fixação $4N = 160$ gerações.
\textbf{11.} Sem endogamia, $q^{2} = 0.0025$; com $F = 0.3$,
$q^{2} + Fpq = 0.0025 + 0.3\times 0.95\times 0.05 = 0.017$: sete
vezes mais aves afetadas.
\textbf{12.} Um migrante por geração ($m = 1/40$) basta para impedir
que a ilha diverja nos lócus neutros e para repor os \omterm{def:b2:meiosis-heredity:vocabulary}{alelos} que a deriva
perdeu; ele também importa os \omterm{def:b2:meiosis-heredity:vocabulary}{alelos} deletérios do continente em sua
frequência continental, os quais, sob a endogamia da ilha, ficam
expostos em \omterm{def:b2:meiosis-heredity:vocabulary}{homozigotos} --- a carga é fixada pelo equilíbrio entre
imigração e exposição.
\textbf{13.} $\hat q = \sqrt{10^{-5}} = 0.0032$; nascimentos afetados $\hat
q^{2} = 10^{-5}$.
\textbf{14.} $\hat q = \sqrt{10^{-5}/0.1} = 0.01$; nascimentos afetados
$10^{-4}$: uma seleção dez vezes mais fraca deixa o \omterm{def:b2:meiosis-heredity:vocabulary}{alelo} ficar três
vezes mais comum e a doença, dez vezes.
\textbf{15.} $\hat p = \mu/(hs) = 10^{-5}/0.05 = 2\times 10^{-4}$;
portadores $2\hat p = 4\times 10^{-4}$.
\textbf{16.} De $\hat q = 0.0032$ a $0.01$: o \omterm{def:b2:meiosis-heredity:vocabulary}{alelo} sobe a uma taxa
fixada pela \omterm{def:b2:genome-diversification:mutation}{mutação}, $\mu = 10^{-5}$ por geração, de modo que a
aproximação leva da ordem de $\Delta q/\mu \approx 700$ gerações ---
\num{20000} anos: qualquer decisão tomada agora é sentida por uma
população que não se parecerá em nada com a nossa.
\textbf{17.} Cada lócus contribui com $2\hat q = 0.0063$ cópias letais
por pessoa; sobre \num{20000} lócus, cerca de $130$ --- mas isso supõe
que todo gene possa mutar para um letal \omterm{def:b2:meiosis-heredity:vocabulary}{recessivo}; com a estimativa
convencional de alguns milhares de genes essenciais, uma pessoa carrega
de vários a algumas dezenas de equivalentes letais em estado \omterm{def:b2:meiosis-heredity:vocabulary}{heterozigoto}.
\textbf{18.} Cada um está num lócus diferente e cada um é raro, de modo
que a chance de um casal ao acaso carregar ambos o mesmo é pequena; primos
compartilham um oitavo de seus genes por descendência, de modo que, para
cada um dos vários letais de um primo, o risco de o filho ser \omterm{def:b2:meiosis-heredity:vocabulary}{homozigoto} é
de $1/16$ por letal compartilhado --- o excesso de mortalidade infantil
dos casamentos entre primos é medido em alguns por cento.
\textbf{19.} $S = 1.16\times \qty{800}{L} = \qty{930}{L}$; $R = 0.4\times
930 = \qty{370}{L}$.
\textbf{20.} Cinco gerações: cerca de \qty{1860}{L}; menos na
prática, porque a variância aditiva se esgota, porque $h^{2}$
foi estimada no rebanho original e porque o ambiente
(alimentação, instalações) que produziu as melhores vacas pode não ser
transmitido.
\textbf{21.} Diferencial médio $(1.16 + 2.4)/2 = 1.78$ desvio-padrão,
$S = \qty{1424}{L}$, $R = 0.4\times 1424 = \qty{570}{L}$.
\textbf{22.} \omterm{def:b2:asexual-reproduction:clone}{Clones} não têm variância genética, de modo que $h^{2} = 0$ e $R =
0$, qualquer que seja $S$: a herdabilidade é uma propriedade de uma
população, a parte de sua variância que é genética, e não uma propriedade
do traço.
\textbf{23.} $R = 0.7\times 0.6 = \qty{0.42}{mm}$ mais profundo.
\textbf{24.} A mesma resposta $\qty{370}{L}$ exige $S = 370/0.1 =
\qty{3700}{L}$, isto é, $4.6$ desvios-padrão --- conservar menos de
uma vaca em dez mil: impraticável, e é por isso que os traços de baixa
herdabilidade são melhorados lentamente ou por outros meios
(teste de progênie, famílias maiores).
\textbf{25.} Mariposa: $s \approx 0.2$ a $0.3$, 28 gerações com $0.3$
e 45 com $0.2$, contra 50 observadas; ilha: $H_{100} \approx 0.11$,
\qty{27}{\%} do continente; rebanho: \qty{370}{L} por geração só com as
vacas, \qty{570}{L} com touros selecionados.
\end{solution}
